\section{The Architect's Question}\label{the-architects-question}

After chapters on tokens, workloads, and memory, the architect naturally
asks: \emph{how much compute does this actually require, and at what
point does the hardware cease to be the limiting factor?} This chapter
gives us the FLOP-scale arithmetic, the arithmetic-intensity roofline,
and the utilization framework that lets us answer that question without
guessing. We will move from per-token FLOPs to prefill PFLOP counts to
sustained MFU on real hardware, and then to the roofline that tells us
whether a given layer is compute-bound or memory-bound --- the
difference that decides whether we should reach for batching,
quantization, or a different hardware generation. Throughout, the
arithmetic anchors to the canonical enterprise-Q\&A RAG scenario (§14):
\textasciitilde10 rps average / \textasciitilde40 rps peak traffic
(2,000 registered users × 5\% concurrency → 100 concurrent →
\textasciitilde10 rps via Little's Law; peaks at 20\% concurrency →
\textasciitilde40 rps), with 70B FP16 weights and an 8×H100 host.
KV-cache per-token figures are referenced at FP16 (\textasciitilde2.5
MB/token) unless an 8-bit (\textasciitilde1.3 MB/token) variant is
explicitly stated.

\section{Concept}\label{concept}

Compute in a language model is priced in floating-point operations, or
FLOPs. Every matrix multiplication, every attention weight update, every
activation function evaluation contributes to the total. The architect's
first principle is that FLOPs scale linearly with parameter count and
context length, but the \emph{efficiency} with which those FLOPs are
executed depends on the hardware ridge point and the arithmetic
intensity of the kernel.

For a dense transformer layer, the dominant cost is the matrix
multiplication --- query-key, value-aggregation, and MLP projections. In
FP16, a single multiply-add counts as two FLOPs. The per-token cost is
therefore approximately 2 × N FLOPs where N is the number of active
parameters. This is the arithmetic baseline: every token processed costs
roughly two floating-point operations per parameter.

But FLOPs alone do not tell the full story. The hardware can only
execute FLOPs if data is available --- weights, activations, and KV
cache all compete for the same HBM bandwidth. The arithmetic intensity
(FLOP/byte) determines whether a kernel is compute-bound (intensity
above the ridge point) or memory-bound (intensity below). This roofline
model is the central diagnostic tool of this chapter: it tells us, for
any given layer and precision, whether increasing FLOPs will actually
reduce latency or whether we are already starved for bytes.

The chapter unfolds in four parts. First, the core arithmetic: FLOPs per
token, prefill PFLOP counts, and H100 sustained utilization. Second, the
roofline: why early layers are compute-bound while later layers and
decode are memory-bound. Third, the canonical scenario arithmetic: 70B
on 8 × H100, TTFT 1.2 s, TPOT \textasciitilde25 ms. Fourth, the
mini-case: a continuous deployment scenario that threads the
prefill/decode divide.

\section{Mental Model}\label{mental-model}

Think of FLOPs as the distance a car can travel on a gallon of fuel: it
tells us the system's capacity, not how long the trip will take. The
trip time is determined by the ridge point: if the arithmetic intensity
is below the ridge, adding more FLOPs won't help because we are waiting
for data. If we are above the ridge, we are limited by how fast the
compute can chew through operations. The roofline chart --- peak TFLOPS
on the vertical axis, FLOP/byte on the horizontal --- is the map that
resolves this. Where our kernel lands on that chart decides whether we
should profile memory bandwidth or compute throughput.

\section{Worked Example: Canonical Scenario
Arithmetic}\label{worked-example-canonical-scenario-arithmetic}

The canonical scenario (§14, book-architecture.md) is an enterprise Q\&A
system: 70B-class dense model, FP16 weights (\textasciitilde140 GB), 1
host with 8 × H100-class GPUs (80 GB each), \textasciitilde10 requests/s
average, peaks \textasciitilde40 rps, average prompt 1,200 tokens + 8K
retrieved context (\textasciitilde9.2K input), 300-token output, TTFT
budget 1.2 s (retrieval \textasciitilde120 ms + prefill), TPOT budget
\textasciitilde25 ms/token.

\subsection{FLOPs per token}\label{flops-per-token}

For a 70B-parameter dense model in FP16, the per-token forward FLOP
count is approximately:

\[
\text{FLOP/}_{\text{token}} \approx 2 \times N = 2 \times 70 \times 10^9 \approx 140 \text{ GFLOP/token}
\]

{[}DERIVED: 2 FLOPs/parameter × 70 × 10⁹ params; standard dense-forward
arithmetic, consistent with the 175B ≈ 0.35 TFLOP/token correction from
moe-vs-dense E5{]}

This applies to both prefill and decode: each token that enters the
forward pass costs \textasciitilde140 GFLOP. The distinction between
prefill and decode is not in the per-token FLOP count but in the data
movement pattern --- preflight reads the full context once, while decode
re-reads weights per token.

\subsection{Prefill PFLOP for 9.2K
input}\label{prefill-pflop-for-9.2k-input}

For a single request with 9.2K input tokens, the prefill FLOP cost is:

\[
\text{prefill FLOPs} \approx 2 \times N \times L = 2 \times 70 \times 10^9 \times 9.2 \times 10^3 \approx 1.29 \text{ PFLOP}
\]

{[}DERIVED: 2 × N × L where N=70B, L=9,200; reconciles with Ch.2 scaling
law if the full 14.8T-token pre-train budget is distributed across
active parameters; the figure is large because the full context is
attended to, but it is a one-time cost per request, not a sustained
rate{]}

\emph{Accuracy of the 2 × N × L approximation.} This linear prefill
count omits the quadratic attention term,
\(\approx 4 \times n_\text{layers} \times L^2 \times d\). At the
canonical 9.2K context that correction is small --- roughly
\textbf{+17\%} of the 1.29 PFLOP --- so the 2NL roofline is a sound
teaching baseline there. But the omission grows with context: at 32K the
attention term is \textasciitilde{}\textbf{+60\%} and at 128K it
\emph{dominates} (\textasciitilde2.4× the linear term). An architect
sizing long-context prefill must add the quadratic term or measure it;
the Unknowns section returns to this.

At 10 requests/s, the sustained prefill throughput demand is:

\[
\text{prefill demand} = 1.29 \times 10^{15} \text{ FLOP} \times 10 \text{ rps} \approx 12.9 \text{ PFLOP/s}
\]

{[}DERIVED{]}

\textbf{\hyperref[tab:8.1]{Table 8-1}} --- Per-token FLOP cost and prefill PFLOP for the
canonical 70B workload. \emph{(Per-token FLOP and prefill PFLOP
computations are {[}2° DERIVED{]} from 2 × params × tokens; hardware
peaks {[}1P: vendor datasheet{]} or {[}2° FACT{]}; the 30--40\% MFU
range is {[}2°: industry benchmarks{]}; the \textasciitilde295 FLOP/byte
ridge is {[}2° DERIVED{]} = 989 TFLOPS ÷ 3.35 TB/s.)} \textbar{} metric
\textbar{} value \textbar{} derivation \textbar{}

\begin{longtable}[]{@{}
  >{\raggedright\arraybackslash}p{0.3333\linewidth}
  >{\raggedright\arraybackslash}p{0.3333\linewidth}
  >{\raggedright\arraybackslash}p{0.3333\linewidth}@{}}
\toprule\noalign{}
\begin{minipage}[b]{\linewidth}\raggedright
metric
\end{minipage} & \begin{minipage}[b]{\linewidth}\raggedright
value
\end{minipage} & \begin{minipage}[b]{\linewidth}\raggedright
derivation
\end{minipage} \\
\midrule\noalign{}
\endhead
\bottomrule\noalign{}
\endlastfoot
per-token FLOPs (forward) & \textasciitilde140 GFLOP/token & 2 × 70B
params {[}DERIVED; consistent with moe-vs-dense E5 correction: 175B ≈
0.35 TFLOP/token{]} \\
prefill FLOPs for 9.2K input & \textasciitilde1.29 PFLOP & 2 × 70 × 10⁹
× 9.2 × 10³ {[}DERIVED{]} \\
sustained prefill demand @ 10 rps & \textasciitilde12.9 PFLOP/s & 1.29
PFLOP × 10 {[}DERIVED{]} \\
H100 peak FP16 TFLOPS & \textasciitilde989 TFLOPS {[}1P:
facts/training.md T5{]} & NVIDIA H100 SXM5 datasheet, without
sparsity \\
H100 sustained MFU (typical) & 30--40\% {[}2°: industry benchmarks{]} &
\textasciitilde346 TFLOPS sustained at 35\% MFU \\
H100 ridge point (dense FP16) & \textasciitilde295 FLOP/byte & 989 ÷
3.35 {[}DERIVED: peak TFLOPS ÷ HBM bandwidth{]} \\

\label{tab:8.1}\end{longtable}

\emph{All figures trace to the canonical scenario (§14) and validated
old-repo sources; none are measurement claims.} This is the prefill
compute demand. An 8 × H100 node can sustain some fraction of this at
MFU (mixed-precision FLOP utilization), which we estimate next.

\subsection{Sustained vs.~peak: H100
MFU}\label{sustained-vs.-peak-h100-mfu}

NVIDIA H100 SXM5 lists \textasciitilde1,979 TFLOPS FP16 Tensor Core peak
``with sparsity'' and \textasciitilde989 TFLOPS FP16 peak without
sparsity {[}1P: facts/training.md T5, NVIDIA datasheet{]}. In practice,
sustained mixed-precision FLOP utilization (MFU) for a well-tuned
transformer workload typically runs at 30--40\% of peak on H100 {[}2°:
industry benchmark reports, derived from real kernel profiles{]}. At
35\% MFU:

\[
\text{sustained} = \text{peak} \times \text{MFU} = 989 \text{ TFLOPS} \times 0.35 \approx 346 \text{ TFLOPS}
\]

{[}DERIVED: peak × typical MFU{]}

At 346 TFLOPS sustained, the number of tokens/s the system can prefill
is:

\[
\text{prefill tokens/s} = \frac{\text{sustained}}{\text{FLOP/}_{\text{token}}} = \frac{346 \text{ TFLOPS}}{140 \text{ GFLOP/token}} \approx 2{,}471 \text{ tokens/s}
\]

{[}DERIVED: sustained FLOPS ÷ per-token FLOPs{]}

For 10 rps with 9.2K input tokens each, the prefill demand is
\textasciitilde92,000 tokens/s (from the canonical workload). At 2,471
tokens/s per node, we would need roughly:

\[
n_\text{nodes} = \frac{\text{tokens/s demand}}{\text{tokens/s per node}} = \frac{92{,}000}{2{,}471} \approx 37 \text{ nodes}
\]

{[}DERIVED: token demand ÷ per-node prefill throughput{]}

This rough sizing illustrates that prefill is FLOP-bound at this scale
--- the compute requirement is the primary constraint, and memory
bandwidth (HBM at 3.35 TB/s per H100) is sufficient to keep the compute
fed if the kernel is efficient. The ridge point for dense FP16 on H100
is \textasciitilde295 FLOP/byte (989 ÷ 3.35); the arithmetic intensity
of a mature attention kernel sits near or above this ridge, meaning the
kernel is compute-saturated rather than bandwidth-starved once the batch
is large enough.

\subsection{Decode: bandwidth-bound at low
batch}\label{decode-bandwidth-bound-at-low-batch}

\begin{figure}
\centering
\pandocbounded{\includegraphics[keepaspectratio,alt={\hyperref[fig:8.1]{Fig 8.1} --- Prefill above the ridge (compute-bound), decode below it (memory-bound). H100 (solid) vs H200 (dashed) ridge points from vendor datasheets (989 TFLOPS, 3.35 / 4.8 TB/s); Continuous-Batching arrow shows decode climbing the slope as batch grows {[}2° DERIVED{]} {[}1P: vendor datasheet{]}},width=\maxwidth{\textwidth},keepaspectratio]{figures/fig-08-0801.pdf}}
\caption{\hyperref[fig:8.1]{Fig 8.1} --- Prefill above the ridge (compute-bound), decode
below it (memory-bound). H100 (solid) vs H200 (dashed) ridge points from
vendor datasheets (989 TFLOPS, 3.35 / 4.8 TB/s); Continuous-Batching
arrow shows decode climbing the slope as batch grows {[}2° DERIVED{]}
{[}1P: vendor datasheet{]}}

\label{fig:8.1}\end{figure}

\emph{\hyperref[fig:8.1]{Fig 8.1} --- The roofline: prefill is compute-bound, low-batch
decode is memory-bound.}

For decode, the per-token FLOP cost is the same \textasciitilde140
GFLOP, but the arithmetic intensity drops sharply. Each decoded token
re-reads all 70B weights from HBM. A 70B FP16 weight matrix is 70 × 10⁹
× 2 B = 140 GB --- one full pass reads \textbf{140 GB}, not 280 GB. (The
naïve ``2 × 70 × 10⁹ × 2 B ≈ 280 GB'' double-counts the FP16 footprint:
the leading ``2'' is already inside the 2-bytes-per-element, so writing
2 × 70 × 10⁹ × 2 B counts the weight bytes twice. Any
KV/output-projection reads are a separate, smaller term on top, not a
second full weight footprint.) The effective arithmetic intensity is
therefore:

\[
\text{arithmetic intensity}_{\text{decode}} \approx \frac{140 \text{ GFLOP}}{140 \text{ GB}} \approx 1.0 \text{ FLOP/byte}
\]

{[}DERIVED; a lower-bound bandwidth model --- real kernels also read
activations and KV, so measured intensity is typically below this and
heavily dependent on fusion and cache residency{]}

This is well below the H100 dense-FP16 ridge of \textasciitilde295
FLOP/byte, meaning single-stream decode is overwhelmingly
memory-bandwidth-bound. Throughput improves with batching because the
weight-reread amortizes across tokens --- at batch=32, the effective
intensity rises toward the ridge, and TFLOP/s utilization increases
accordingly. This is why continuous-batching systems (Orca, vLLM)
prioritize batch growth: it is the primary lever for pushing decode
arithmetic intensity toward the ridge.

\section{Measurement}\label{measurement}

Three measurable quantities anchor compute sizing:

\begin{enumerate}
\def\labelenumi{\arabic{enumi}.}
\item
  \textbf{Per-token FLOP count.} Measure by flops profiling (e.g.~NVIDIA
  Nsight Compute) on the target model and precision. The 2×N rule is a
  useful prior, but actual count varies by layer depth, activation
  choice (SiLU vs.~ReLU), and whether kernel fusion combines matmuls.
\item
  \textbf{Sustained MFU.} Profile TFLOP/s achieved divided by peak
  TFLOP/s for the target hardware and precision. A well-tuned
  transformer on H100 FP16 typically lands in the 30--40\% range; below
  20\% suggests a memory or communication bottleneck that should be
  diagnosed before scaling.
\item
  \textbf{Arithmetic intensity (FLOP/byte).} Compute as achieved FLOPs
  divided by bytes transferred (weights + activations + KV read/write).
  Compare to the hardware ridge point (H100 dense FP16:
  \textasciitilde295 FLOP/byte; H200: \textasciitilde989 TFLOPS ÷ 4.8
  TB/s ≈ 206 FLOP/byte --- H200 keeps H100's GH100 die and its dense
  FP16 peak, trading a bigger/faster HBM3e for bandwidth; GB200 BF16:
  much higher due to NVFP4). If intensity is below the ridge, the kernel
  is memory-bound; above, compute-bound.
\end{enumerate}

\section{Common Mistakes}\label{common-mistakes}

\begin{itemize}
\item
  \textbf{Assuming 2×N FLOPs/token is exact.} It is a baseline
  derivation; actual per-token FLOPs vary with layer depth, activation
  fusion, and kernel fusion. Profile before quoting.
\item
  \textbf{Treating prefill and decode as having the same FLOP/byte
  profile.} Prefill reads the full context once and is FLOP-bound;
  decode re-reads weights per token and is memory-bound. The
  arithmetic-intensity roofline distinguishes them.
\item
  \textbf{Using peak TFLOP/s as sustained throughput.} Peak includes
  sparsity or BF16/FP8 variants that may not be available in the chosen
  precision. Use the FP16 number without sparsity (≈989 TFLOPS on H100)
  as the baseline for dense compute sizing.
\item
  \textbf{Ignoring the batch-size dependence of arithmetic intensity.} A
  single-token decode batch has very low intensity; batching raises it.
  Always state the batch assumption when reporting TFLOP/s or tokens/s.
\item
  \textbf{Confusing H100 sparsity-inclusive peak (1,979 TFLOPS) with
  dense peak (989 TFLOPS).} The 2× sparsity figure appears in vendor
  data sheets but most transformer workloads do not activate sparsity.
  Quote the dense figure unless the kernel is specifically
  sparsity-enabled.
\end{itemize}

\section{Architecture Consequence}\label{architecture-consequence}

The compute arithmetic and roofline have direct consequences for system
design:

\begin{itemize}
\item
  \textbf{If prefill is FLOP-bound} (as it is for 70B+ models with long
  contexts), the primary optimization is batching --- larger prefill
  batches amortize the context attention cost across more requests,
  raising arithmetic intensity and increasing TFLOP/s utilization.
  System design should therefore provision for batchable prefill
  (continuous batching, Orca, vLLM continuous batching) as the
  first-order throughput lever.
\item
  \textbf{If decode is memory-bound} (single-stream, low batch), the
  primary optimization is batching --- even modest batch sizes (8--16)
  raise arithmetic intensity toward the ridge, improving TFLOP/s
  utilization. KV cache quantization (FP8, int8) also reduces bytes per
  token, raising effective intensity. P/D disaggregation (prefill on
  GPUs, decode on a separate pool) can rebalance: prefill's FLOP demand
  is served by compute-optimized GPUs, while decode's bandwidth demand
  is served by wide-memory GPUs or even CPU offload.
\item
  \textbf{If the ridge point is exceeded} (e.g.~H200 with 4.8 TB/s HBM3e
  and 4 PFLOPS FP8, giving a \textasciitilde833 FLOP/byte ridge in FP8),
  the same workload may shift from memory-bound to compute-bound,
  changing the optimal hardware choice. This is why the H200 represents
  a inflection point where inference and training hardware lines begin
  to converge {[}1P: facts/training.md T5{]}.
\end{itemize}

In practice, the architect measures arithmetic intensity for the target
workload and precision, locates the kernel on the roofline chart, and
then selects hardware and batching strategy accordingly. The roofline is
the diagnostic; the architecture decision follows.

\section{What We Still Don't Know}\label{what-we-still-dont-know}

\begin{itemize}
\item
  \textbf{Arithmetic intensity of fused kernels.} Most roofline
  estimates treat attention or MLP as separate matmuls, but in practice
  frameworks fuse multiple operations into a single kernel. The
  effective FLOP/byte of a fused attention+MLP kernel is not well
  cataloged in primary sources and would require profiling on the target
  hardware to anchor.
\item
  \textbf{MFU across precisions and sparsity.} The 30--40\% H100 FP16
  figure is a rule of thumb; actual utilization depends on architecture,
  kernel quality, and batch config. A systematic MFU atlas across
  H100/H200/GB200, FP16/BF16/FP8, dense/MOE would be a valuable
  reference.
\item
  \textbf{Arithmetic intensity of emerging attention mechanisms.} Linear
  attention (DeltaNet, MLA), compressed attention (CSA, HCA), and hybrid
  sparse patterns have different FLOP profiles and data movement. Their
  roofline position is not yet established in primary sources.
\item
  \textbf{Frontier 2026 has begun to answer this.} DeepSeek-V4's hybrid
  CSA+HCA attention reports that, at 1M-token context, single-token
  inference drops to \textasciitilde27\% of the FLOPs (Pro) /
  \textasciitilde10\% (Flash) and KV cache to \textasciitilde10\% /
  \textasciitilde7\% of DeepSeek-V3.2 {[}1P: arXiv 2606.19348{]}.
  GLM-5.3-Flash's sparse+linear hybrid reports a \textasciitilde3×
  attention-compute and \textasciitilde4.4× KV-cache reduction {[}1P: HF
  zai-org/GLM-5.3-Flash{]}. These are {[}1P{]} first-party vendor
  figures, not yet independently reprofiled on our own hardware ---
  which is exactly the arithmetic-intensity measurement an architect
  should still do before trusting a vendor's roofline claim for their
  own workload.
\item
  \textbf{How quantization shifts the ridge.} Moving from FP16 to
  FP8/int8 changes peak TFLOP/s and bytes-per-operand, shifting the
  ridge point. The net effect depends on the quantization scheme
  (element-wise vs.~per-tensor, dynamic vs.~static) and whether the
  kernel is re-tuned for the lower precision.
\end{itemize}

\section{End-of-Chapter Mini-Case: Continuous Deployment
Scenario}\label{end-of-chapter-mini-case-continuous-deployment-scenario}

An architect is brought into an ongoing deployment of a 70B-class Q\&A
system on 8 × H100 GPUs. The system is serving \textasciitilde10
requests/s average with \textasciitilde9.2K input + 300 output tokens
per request, and the TTFT budget is being missed: p95 TTFT is 2.8 s,
exceeding the SLO of 2 s. The TPOT of 28 ms/token is within spec, but
the prefill delay is the bottleneck.

The architect first profiles the arithmetic intensity of the prefill
kernel. The per-token FLOP count is \textasciitilde140 GFLOP (2 × 70B),
and the effective bandwidth per token is \textasciitilde140 GB
(re-reading the 70B weight matrix from HBM for each of the 9.2K input
tokens). This gives an effective FLOP/byte of roughly 140 GFLOP ÷ 140 GB
≈ 1 FLOP/byte --- well below the H100 dense-FP16 ridge of
\textasciitilde295 FLOP/byte. The kernel is strongly memory-bound during
prefill.

The immediate question: why is prefill memory-bound when the H100 has
3.35 TB/s HBM3 bandwidth? The answer is that a 9.2K input context means
each prefill pass reads the full weight matrix once, but the activation
keys/values for 9.2K tokens also flow through the pipeline. The total
bytes per token are higher than the weight-read alone, and the effective
intensity drops further.

The architect's first fix is to increase the prefill batch size.
Continuous batching (vLLM Orca-style) accumulates incoming requests into
larger batches, which raises the effective arithmetic intensity because
the weight matrix read is amortized across more tokens. At batch=32, the
effective intensity rises to \textasciitilde2--3 FLOP/byte, still below
the ridge but much closer. At batch=128, the intensity approaches the
ridge and TFLOP/s utilization climbs from \textasciitilde15\% to
\textasciitilde45\%.

The second fix is to profile the actual MFU achieved. If the system is
running at 12\% MFU on prefill, there is headroom to absorb more batch
without adding hardware. If MFU is already near 40\%, the bottleneck is
elsewhere (communication, kernel inefficiency, or KV cache residency),
and simply batching harder yields diminishing returns.

The architect's recommendation: enable continuous batching with a target
batch size of 64--128, profile the resulting MFU and TTFT, and if TTFT
remains above SLO, consider P/D disaggregation --- offload decode to a
separate pool of GPUs with wide HBM bandwidth, freeing the prefill GPUs
to run larger batches at higher FLOP utilization. The key insight is
that the workload's position on the roofline chart (memory-bound
prefill) dictates the fix: batching to raise intensity, or hardware that
raises the ridge point (H200, GB200), or disaggregation that separates
the two bottleneck profiles.

\begin{center}\rule{0.5\linewidth}{0.5pt}\end{center}

\emph{Reading the roofline (summary of \hyperref[fig:8.1]{Fig 8.1}): the dense-FP16 ridge at
\textasciitilde295 FLOP/byte (H100: 989 TFLOPS ÷ 3.35 TB/s) separates
the compute-bound region (right) from the memory-bound region (left).
Prefill kernels typically operate in the 1--3 FLOP/byte range for long
contexts, decode kernels near 1.0 FLOP/byte at batch=1; batching and
quantization shift kernels rightward toward the ridge.}

\emph{Prefill PFLOP also grows with context: 2 × 70B × L tokens, from 2K
to 32K. This is the one-time compute cost of loading a long context;
sustained throughput (tokens/s) is what matters for continuous serving,
and at the canonical 9.2K point it is \textasciitilde1.29 PFLOP (\hyperref[tab:8.1]{Table 8-1}).}
